% !TEX program = xelatex
% ---------------------------------------------------------------------------
%  胡霄星 · 学术简历（中文版）
%  编译：xelatex CV-zh.tex
% ---------------------------------------------------------------------------
\documentclass[a4paper,11pt]{ctexart}

\usepackage[left=1.75cm,right=1.75cm,top=1.6cm,bottom=1.6cm]{geometry}
\usepackage{xcolor}
\usepackage{titlesec}
\usepackage{enumitem}
\usepackage{tabularx}
\usepackage{array}
\usepackage{fontspec}
\usepackage{xeCJK}
\usepackage{needspace}
\usepackage{fontawesome5}
\usepackage[hidelinks]{hyperref}
\hypersetup{linkcolor=black, urlcolor=black, citecolor=black}

% ── 配色 ────────────────────────────────────────────────────────────────────
\definecolor{accent}{HTML}{000000}
\definecolor{muted}{HTML}{000000}
\definecolor{rule}{HTML}{000000}

% ── 字体 ────────────────────────────────────────────────────────────────────
\setCJKmainfont{Kaiti SC}[BoldFont=Songti SC Bold]
\newCJKfontfamily\xingkai{行楷-简}[BoldFont=行楷-简 粗体]
\newCJKfontfamily\kaiti{Kaiti SC}[BoldFont=Kaiti SC Bold]
\newCJKfontfamily\xhei{Heiti SC}[BoldFont=Heiti SC, ItalicFont=Heiti SC,
  BoldFeatures={FakeBold=1.8}, ItalicFeatures={FakeSlant=0.2}]
\setmainfont{Charter}[
  BoldFont       = Charter Bold,
  ItalicFont     = Charter Italic,
  BoldItalicFont = Charter Bold Italic
]
\setsansfont{Helvetica Neue}
\setmonofont{Menlo}

\pagestyle{empty}
\pagecolor{white}
\setlength{\parindent}{0pt}
\linespread{1.22}
\setlength{\parskip}{0pt}

% ── 标题样式 ────────────────────────────────────────────────────────────────
\titleformat{\section}
  {\large\bfseries\color{accent}}
  {}{0em}{}
  [\vspace{-0.85em}{\color{rule}\rule{\textwidth}{0.5pt}}]
\titlespacing*{\section}{0pt}{1.1em}{0.55em}
% keep a heading with the content that follows it
\let\oldSection\section
\renewcommand{\section}[1]{\needspace{5\baselineskip}\oldSection{#1}}

% ── 列表 ────────────────────────────────────────────────────────────────────
\setlist[itemize]{
  leftmargin=1.35em,
  topsep=2pt,
  itemsep=1.5pt,
  parsep=0pt,
  label={\color{accent}\footnotesize\textbullet}
}

% ── 自定义命令 ──────────────────────────────────────────────────────────────
% \entry{机构}{角色}{时间}{可选说明}
\newcommand{\entry}[4]{%
  \begin{tabularx}{\textwidth}{@{}X r@{}}
    {\bfseries #1}\hfill{\kaiti\itshape\color{muted}#2} & {\kaiti\small\color{muted}#3}
  \end{tabularx}%
  \par\vspace{1.5pt}%
  \if\relax\detokenize{#4}\relax\else{\kaiti\small #4}\par\fi
  \vspace{0.55em}%
}

% \pub{标题}{作者}{会议}{链接}
\newcommand{\pub}[4]{%
  \needspace{4\baselineskip}%
  \vspace{0.32em}%
  {\bfseries #1}\par\vspace{0.5pt}%
  {\kaiti\small #2}\par\vspace{0.5pt}%
  {\kaiti\small\itshape\color{accent}#3}\par\vspace{0.5pt}%
  {\kaiti\footnotesize\color{muted}#4}\par
}

\newcommand{\sep}{\,\textperiodcentered\,}

% ---------------------------------------------------------------------------
\begin{document}

% ── 抬头 ────────────────────────────────────────────────────────────────────
\begin{center}
  {\xingkai\fontsize{34}{40}\selectfont 胡霄星}\\[8pt]
  {\small
    \faEnvelope\ \href{mailto:xiaoxinghhh@gmail.com}{xiaoxinghhh@gmail.com}\quad
    \faPhone\ 15111480307\quad
    \faWeixin\ 15111480307\quad
    \faGoogle\ \href{https://scholar.google.com/citations?user=zBM8_XkAAAAJ&hl=zh-CN}{Google Scholar}
  }
\end{center}

\vspace{-0.35em}
{\color{rule}\rule{\textwidth}{0.6pt}}

% ── 教育经历 ────────────────────────────────────────────────────────────────
\section{教育经历}

\entry{上海交通大学}{博士研究生，控制科学与工程}{2026.09 -- 至今}
  {}

\entry{上海创智学院}{博士研究生（联合培养）}{2026.09 -- 至今}
  {}

% ── 研究兴趣 ────────────────────────────────────────────────────────────────
\section{研究兴趣}

{\xhei\bfseries 研究经历：}过往研究聚焦于视觉基础模型及其遥感交叉应用、跨模态表征与对齐，以及（音）视频理解与评测。具体工作涵盖视觉--语言对齐与多模态检索、证据可溯源的多模态推理基准与自动化评测，以及视觉基础模型的迁移学习与域泛化。

\vspace{0.45em}
{\xhei\bfseries 当前研究兴趣：}聚焦于智能体全栈技术，涵盖模型、数据、系统架构及大语言模型后训练（OPD、MOPD、SWE-OPD）等方向。

% ── 论文发表 ────────────────────────────────────────────────────────────────
\section{论文发表}
{\xhei\bfseries\footnotesize\color{muted}加粗为本人。\quad
 $^{*}$ 共同一作。\quad $^{\dagger}$ 通讯作者。\quad $^{\ddagger}$ 学生一作。\par}
\vspace{0.2em}

\pub{EvaOmni: An Evidence-Grounded Benchmark for Audio-Visual Reasoning}
  {\textbf{Xiaoxing Hu}, Wenxi Yao, Chaoyuan Shi, Ziqi Ye, Xiaorong Guo, Peixuan Li, Zhuoran Sun,
   Ziyang Gong, Peiyuan Zhang, Mingxin Liu, Shuran Ma, Hongbo Liu, Xiaomin Yu, Yu Tian, Wentao Wang,
   Zhenjie Yang, Jian Hu, Yunhang Shen, Haoyu Cao, Xue Yang$^{\dagger}$}
  {审稿中}
  {\href{https://xiaoxing2001.github.io/EvaOmni.github.io/}{项目主页}\quad
   \href{https://github.com/VisionXLab/EvaOmni}{代码}\quad
   \href{https://arxiv.org/abs/2610.28444}{论文}}

\pub{ProCLIP: Progressive Vision-Language Alignment via LLM-based Embedder}
  {\textbf{Xiaoxing Hu}, Kaicheng Yang, Ziyang Gong, Qi Ming, Zonghao Guo, Yu Tian, Xiang An,
   Ziyong Feng, Xue Yang$^{\dagger}$}
  {NeurIPS 2026 (Spotlight)\sep CCF-A}
  {\href{https://github.com/VisionXLab/ProCLIP}{代码}\quad\href{https://arxiv.org/abs/2510.18795}{论文}}

\pub{Earth-Adapter: Bridge the Geospatial Domain Gaps with Mixture of Frequency Adaptation}
  {\textbf{Xiaoxing Hu}$^{*}$, Ziyang Gong$^{*}$, Yupei Wang$^{\dagger}$, Yuru Jia, Fei Lin, Dexiang Gao, Ke An,
   Jianhong Han, Zhuoran Sun, Gen Luo, Xue Yang$^{\dagger}$}
  {AAAI 2026\sep CCF-A}
  {\href{https://github.com/VisionXLab/Earth-Adapter}{代码}\quad
   \href{https://arxiv.org/abs/2504.06220}{论文}}

\pub{Decoupled Global-Local Alignment for Improving Compositional Understanding}
  {\textbf{Xiaoxing Hu}$^{*}$, Kaicheng Yang$^{*}$, Jun Wang, Haoran Xu, Ziyong Feng, Yupei Wang$^{\dagger}$}
  {ACM MM 2025\sep CCF-A}
  {\href{https://xiaoxing2001.github.io/DeGLA.github.io/}{项目主页}\quad
   \href{https://github.com/xiaoxing2001/DeGLA}{代码}\quad
   \href{https://arxiv.org/abs/2504.16801}{论文}}

\pub{Boosting Domain Generalization in Remote Sensing Image Segmentation via Style Mapping and General Prototypical Contrast}
  {Yupei Wang, \textbf{Xiaoxing Hu}$^{\ddagger}$, Yongkang Hu, Zhuoran Sun, Shanghang Zhang, Liang Chen$^{\dagger}$}
  {International Journal of Computer Vision (IJCV) 2025\sep CCF-A}
  {\href{https://link.springer.com/article/10.1007/s11263-025-02568-3}{论文}}

% ── 实习经历 ────────────────────────────────────────────────────────────────
\section{实习经历}

\entry{腾讯优图\sep 神农研究中心}{研究实习生}{2025.12 -- 2026.09}
  {\begin{itemize}
     \item 参与全模态模型研发，负责数据、训练与评测链路的搭建与迭代。
     \item 参与办公智能体评测体系搭建，针对 PPT-Agent 执行轨迹引入基于 Rubric 的 LLM-as-a-Judge 自动化评测机制，支撑智能体能力评估与迭代验证。
     \item 参与数据集评测与清洗流程建设，撰写并投稿论文一篇（审稿中）。
   \end{itemize}}

\entry{格灵深瞳\sep 大模型组}{研究实习生}{2025.06 -- 2025.10}
  {\begin{itemize}
     \item 负责多模态检索基础模型研究，围绕视觉--语言表征对齐展开实验与分析。
     \item 相关成果以第一作者投稿，并被 NeurIPS 2026 接收为 Spotlight（ProCLIP）。
   \end{itemize}}

% ── 荣誉与奖励 ──────────────────────────────────────────────────────────────
\section{荣誉与奖励}

\begin{itemize}
  \item {\kaiti\bfseries 徐特立奖学金候选}——全校 6 名研究生入选\hfill {\small 2026.06}
  \item {\kaiti\bfseries 国家奖学金}\hfill {\small 2025.10}
\end{itemize}

% ── 技能 ────────────────────────────────────────────────────────────────────
\section{技能}

\begin{itemize}
  \item \textbf{编程语言}：Python（熟练）、C/C++、Shell
  \item \textbf{工具}：PyTorch、Hugging Face Transformers、DeepSpeed、Git、\LaTeX
  \item \textbf{方向技能}：多模态预训练与对齐、指令微调、检索、评测基准构建
\end{itemize}

\end{document}
